\documentclass[11pt]{article}
\usepackage[margin=1in]{geometry}
\usepackage{amsmath}
\usepackage{amssymb}
\usepackage{booktabs}
\usepackage{array}
\usepackage{hyperref}

\title{Unified Picker: Multiplicative Score Formula\\
       \large Regime A (default), with Regime B extension}
\author{sharpsat-separator (Phase 1)}
\date{}

\begin{document}
\maketitle

\section{Score formula (Regime A, $\gamma = 0$)}

For each active branch candidate $x$ in the current sub-component, the picker
computes
\[
S(x) \;=\; \operatorname{base}(x) \cdot \operatorname{boost}(x),
\qquad
\operatorname{pick} \;=\; \arg\max_{x \in \text{candidates}} S(x).
\]

\subsection*{Type-pure base score}

\[
\operatorname{base}(x) \;=\;
\begin{cases}
\, w_v \cdot \max\!\bigl(\varepsilon,\; \operatorname{freq}(v) + 10\,\operatorname{activity}(v) + w_c \cdot \operatorname{cheap}(v)\bigr)
   & x = v \text{ (variable)} \\[4pt]
\, w_C \cdot \sigma\!\bigl(\beta\,(L(C) - m)\bigr)
   & x = C \text{ (clause)}
\end{cases}
\]


The clause-base parameters $\beta$ and $m$ control the shape of the
length-sigmoid.

\noindent The length-decayed clause-density proxy:
\[
\operatorname{cheap}(v) \;=\; \sum_{C \,\ni\, v,\;\text{active}} 2^{-\alpha_d \, |C|_{\text{active}}}
\;+\; \operatorname{binPart}(v) \cdot 2^{-2\alpha_d},
\]
where
$\text{binPart}(v)$
is the count of length-2 clauses still actively constraining $v$
under the current partial assignment.

\subsection*{Multiplicative separator boost}

\[
\operatorname{boost}(x) \;=\; 1 \;+\; \alpha \cdot \exp(-\lambda \cdot \operatorname{rel}_k) \cdot \mathbf{1}[x \in \mathrm{sep}],
\]

\[
\operatorname{rel}_k \;=\; \frac{\#\{\text{active separator elements}\}}{\#\{\text{active variables in comp}\}}.
\]

Above we use the bold-one indicator
\[
\mathbf{1}[P] \;=\; \begin{cases} 1 & \text{if predicate } P \text{ is true} \\ 0 & \text{otherwise.} \end{cases}
\]

\subsection*{Properties}

\begin{itemize}
  \item $\operatorname{base}(x) > 0$ for every active candidate (the
        $\varepsilon$ floor on var-base prevents the degenerate
        all-zero case), hence $S(x) > 0$ always. No ``no candidate found'' edge case.
  \item $\operatorname{boost}(x) \in [1,\; 1+\alpha]$ --- never penalizes a candidate;
        only multiplicatively rewards separator membership.
  \item Within a type, ranking is monotone in $\operatorname{base}$: at a
        given node, $\operatorname{boost}$ is the same across all candidates of one
        type that share the same separator-membership status.
  \item Cross-type competition (var vs.\ clause) happens entirely via the
        ratio $w_v \, \overline{\operatorname{base}_v} \;:\; w_C \, \overline{\operatorname{base}_C}$.
        The defaults $w_v = 1, w_C = 7$ are calibrated so a length-4 clause
        ($\sigma(1) \approx 0.731$) has base $\approx 5$, comparable to a typical
        active variable's $\operatorname{freq} + 20\,\operatorname{activity}$.
\end{itemize}

\section{Regime B extension (not yet enabled; $\gamma > 0$)}

When the cascade boost is enabled, the multiplier acquires a third additive
component proportional to how strongly the variable's BCP-cascade dominates
its peers in the current sub-component:

\[
\operatorname{boost}_{\text{Regime B}}(v) \;=\; 1
   \;+\; \alpha \cdot e^{-\lambda \, \operatorname{rel}_k} \cdot \mathbf{1}[v \in \mathrm{sep}]
   \;+\; \gamma \cdot \max\!\Bigl(0,\; \frac{\operatorname{depth}(v)}{\operatorname{depth}_{\mathrm{med}}} \;-\; 1\Bigr),
\]

\[
\operatorname{depth}(v) \;=\; \log_2 \max\bigl(1,\; \mathtt{computeCascadeScore}(v)\bigr),
\]

with $\operatorname{depth}_{\mathrm{med}}$ estimated as the median over a uniform
random sample of $K = 30$ active variables in the comp (full-set median is
prohibitively expensive per pickBranchTarget call).

The boost is \emph{relative to the local median}: vars at or below the median
contribute $0$, and the contribution grows linearly above it. A var with
$\operatorname{depth}(v) = 2 \cdot \operatorname{depth}_{\mathrm{med}}$ adds $\gamma$;
a var with $5 \cdot \operatorname{depth}_{\mathrm{med}}$ adds $4\gamma$.

This calibrates automatically per-node: cascade-rich components reward only
top-cascade outliers; cascade-poor components recognize modest cascades as
locally exceptional. No absolute threshold.

For clauses, $\operatorname{boost}_{\text{cas}}(C) = 0$ --- the cascade boost
applies only to variables.

\section{Symbol table}

\begin{center}
\begin{tabular}{lll>{\raggedright\arraybackslash}p{8cm}}
\toprule
Symbol & Code knob & Default & Role \\
\midrule
$w_v$         & \texttt{picker\_var\_weight}        & $1.0$   & Variable-base anchor. \\
$w_C$         & \texttt{picker\_clause\_weight}     & $7.0$   & Clause-base anchor (calibrated against typical var-base on real instances; see \S\ref{sec:cal}). \\
$w_c$         & \texttt{cheap\_score\_weight}       & $0.0$   & Cheap-score addend weight inside variable base; default off. \\
$\varepsilon$ & \texttt{picker\_base\_floor}        & $0.01$  & Floor on variable base, keeps $S(x) > 0$. \\
$\beta$       & \texttt{clause\_length\_steepness}  & $1.0$   & Sigmoid slope in clause base. \\
$m$           & \texttt{clause\_length\_midpoint}   & $3.0$   & Sigmoid midpoint in clause base. \\
$\alpha$      & \texttt{picker\_alpha}              & $15.0$  & Separator-boost gain; sep candidates at $\operatorname{rel}_k = 0$ get a $1+\alpha = 16\times$ multiplier. \\
$\lambda$     & \texttt{picker\_lambda}             & $5.0$   & Separator-boost decay rate; smooth fade as separator grows relative to comp. \\
$\gamma$      & \texttt{picker\_gamma}              & $0.0$   & Cascade-boost gain; $0$ in Regime A. \\
$\alpha_d$    & \texttt{stage0\_length\_decay}      & auto    & Cheap-score length decay; auto-tuned at preprocess. \\
\bottomrule
\end{tabular}
\end{center}

\section{Definitions}

\noindent The variable-base components are defined as follows.

\paragraph{Frequency.}
$\operatorname{freq}(v)$ is the number of active original clauses containing
the variable $v$ (either polarity), counted at the current node:
\[
\operatorname{freq}(v) \;=\; \#\bigl\{\, C \in \text{active original clauses}
\;:\; v \in \operatorname{vars}(C) \,\bigr\}.
\]
``Active'' here means the clause is not yet satisfied or removed under the
current partial assignment, and contains at least one unassigned literal.
This is the length-agnostic clause-count proxy for $v$'s structural
prominence.

\paragraph{Activity (VSIDS).}
$\operatorname{activity}(v)$ is the sum over both polarities of the
recency-weighted ``conflict involvement'' counter from the
\emph{Variable State Independent Decaying Sum} heuristic:
\[
\operatorname{activity}(v) \;=\; a^{+}(v) \;+\; a^{-}(v),
\]
where $a^{+}(v)$ and $a^{-}(v)$ are the per-literal activity scores
(\texttt{literal(LiteralID(v,sign)).activity\_score\_} in the code) maintained
as follows during search:
\begin{itemize}
	\item All activities start at $0$.
	\item On every conflict, each literal that participated in the conflict
	analysis (UIP-derived learned clause, etc.) has its activity
	\emph{bumped} by an increment $\Delta_{\mathrm{bump}}$.
	\item Every $N$ decisions (in this codebase, every $128$), all activities
	are scaled by a decay factor $0 < d < 1$
	(\texttt{decayActivities()} in \texttt{solver.cpp}).
\end{itemize}
This produces a recency-biased measure of how often the literal has
recently appeared in conflicts. The factor $10$ in the variable-base
formula scales the summed two-polarity activity into the same magnitude
range as $\operatorname{freq}$ ($\sim 0\text{--}100$ in typical search,
versus $\operatorname{freq} \in 0\text{--}10$); empirical default in
\texttt{Solver::scoreOf}.


\subsection*{Clause length weighting for clause branching}

\paragraph{Binary clauses part}

\[
\operatorname{binPart}(v) \;=\;
\#\bigl\{ C = (\ell_v, \ell_b) \in \mathcal{B}
\;:\; \ell_v \text{ is a literal of } v,\;
\ell_v \text{ unassigned},\;
\ell_b \text{ unassigned} \bigr\}
\]
is the count of \emph{live} (currently neither satisfied nor unit) binary
clauses in $\mathcal{B}$ containing $v$. Here $\mathcal{B}$ is the set of
binary clauses stored in \texttt{binary\_links\_} --- original binaries plus
in-scope learned binaries.
Equivalently: the number of length-2 clauses still actively constraining $v$
under the current partial assignment.

\paragraph{The standard sigmoid}
\[
\sigma(z) \;=\; \frac{1}{1 + e^{-z}},
\]


\begin{itemize}
	\item $m$ (\emph{midpoint}, default $3$) is the active clause length at which
	the sigmoid output equals $\tfrac{1}{2}$. Below $m$, the base shrinks
	toward $0$; above $m$, it saturates toward $1$.
	\item $\beta$ (\emph{steepness}, default $1$) controls how sharply the
	transition occurs around $m$. Small $\beta$ ($\to 0$) flattens the
	curve and treats all clause lengths almost equally; large $\beta$
	($\to \infty$) approaches a hard step at $L = m$. Concrete values at
	$m = 3$:
	\begin{center}
		\begin{tabular}{ccccc}
			\toprule
			$L$ & $\beta = 0.3$ & $\beta = 1$ (default) & $\beta = 5$ \\
			\midrule
			$2$ & $0.43$ & $0.27$ & $0.007$ \\
			$3$ & $0.50$ & $0.50$ & $0.50$  \\
			$4$ & $0.57$ & $0.73$ & $0.99$  \\
			$5$ & $0.65$ & $0.88$ & $\approx 1$ \\
			\bottomrule
		\end{tabular}
	\end{center}
	Default $\beta = 1$ gives a moderate gradient: length-$4$ scores
	$\sim 1.5\times$ length-$3$, length-$5$ scores $\sim 1.8\times$ length-$3$.
	This implements a \emph{soft} preference for longer clauses
	(longer $\Rightarrow$ more lits pinned in the negate-all branch
	$\Rightarrow$ deeper BCP cascade), without forcing a hard threshold.
\end{itemize}


\section{Calibration rationale}\label{sec:cal}

The default $w_C = 7$ is not arbitrary. On instances with $\operatorname{freq}(v) \approx 3\text{--}5$
(typical for sparse 3-SAT with $\sim$1:1 var:clause ratio) and $\operatorname{activity}(v) \approx 0$
(early in search, before learning has fired), the variable base is in $[3, 5]$. A length-4 clause
has $\sigma(\beta(4 - 3)) = \sigma(1) \approx 0.731$. Setting $w_C \cdot 0.731 \approx 5$
gives $w_C \approx 7$, equating clause base to typical var base for $L = 4$ clauses.

The default $\alpha = 15$ was chosen so that a separator candidate at $\operatorname{rel}_k = 0$
gets a $16\times$ multiplier on its base, dominating any non-separator candidate by a wide
margin. This approximates the legacy \texttt{-sepVarBias} mode's near-strict ``consume
separator first'' preference.

These defaults are starting points for the manual sensitivity sweep prescribed in
docs/unified\_picker\_redesign.md (Phase 4). The doc's originally-suggested ranges
(e.g.\ $w_C \in [0.3, 3.0]$) were two orders of magnitude too small relative to the
sigmoid-clamped clause-base range $(0, 1)$ once real var-base magnitudes were measured;
the calibration here resolves that.

\section{Worked example (Regime A)}

A sub-component with $N_{\text{active vars}} = 80$, separator size $3$ ($2$ vars $+ 1$ clause):
\[
\operatorname{rel}_k = \tfrac{3}{80} = 0.0375, \qquad
\alpha \cdot e^{-\lambda \cdot 0.0375} = 15 \cdot e^{-0.1875} \approx 15 \cdot 0.83 \approx 12.4.
\]
A separator candidate gets a multiplier $1 + 12.4 = 13.4\times$ on its base. With
default weights, a length-3 separator clause has base $w_C \cdot \sigma(0) = 7 \cdot 0.5 = 3.5$,
yielding $S = 3.5 \cdot 13.4 \approx 47$. A typical non-separator var with base $\approx 5$
yields $S = 5$. The separator clause wins decisively.

When the separator inflates to $20$ elements ($\operatorname{rel}_k = 0.25$):
\[
\alpha \cdot e^{-1.25} \approx 15 \cdot 0.287 \approx 4.3,
\]
the multiplier drops to $5.3\times$, and the separator clause's $S \approx 18.6$ now competes
much more closely with high-base non-separator candidates --- the smooth handover regime the
formula was designed to produce, with no thresholds.

\end{document}
